\documentclass{article}
\usepackage{graphicx} % Required for inserting images
\usepackage{amsmath}
\usepackage{amssymb}
\usepackage{comment}
\usepackage{amsthm}
\usepackage{color}
%\setlength\parindent{0pt}


\newtheorem*{remark}{Remark}
\newtheorem*{problem}{Problem}
\newtheorem*{theorem}{Theorem}
\newtheorem*{proposition}{Proposition}


\title{Numerical approximation of an optimization problem that appears in  data-driven Computational Mechanics}

\begin{document}

\maketitle

\section{The model problem}

We consider a reaction-diffusion problem, satisfying the conservation law
\begin{equation}
\nabla \cdot s + \zeta\,u = q
\label{eq:conservation}
\end{equation}
in a regular domain $\Omega \subset \mathbb{R}^d$. In (\ref{eq:conservation}) the symbol $\nabla$ represents the vector of partial derivatives, so that $\nabla \cdot s$ is the divergence of the vector field $s$, which represents the flux. Also, $u$ is the concentration of the conserved quantity, $\zeta\in L^{\infty}(\Omega)$ is the (possibly position-dependent) reaction coefficient and $q\in H^1(\Omega)'$ is the source. In thermal applications $u$ is the temperature, $s$ is the heat flux and $q-\zeta u$ is a heat source that depends linearly on the temperature. In weakly compressible flow in porous media $u$ is the pressure and $s$ is the flux.

The boundary of $\Omega$, denoted by $\partial \Omega$, is divided into a Dirichlet part $\partial \Omega_D$, where $u=u_D$ is imposed, and a Neumann part $\partial\Omega_N$, in which $s\cdot \mathbf{\check{n}}=H$ is imposed. Hereafter both $u_D$ and $H$ will be taken as zero to simplify the exposition. 

Let us also define the field $e$ as the gradient of $u$,
\begin{equation}
e = \nabla u~.
\label{eq:compatibility}
\end{equation}
This equation can be viewed as a compatibility condition. We call a vector field $e$ compatible if there exists a scalar field $u$ such that $e=\nabla u$.

Assume now that two vector fields $\tilde{s}$ and $\tilde{e}$ are given in $\Omega$. These fields will both be assumed in $L^2(\Omega)^d$ and thus possibly discontinuous. They come from some measured dataset and thus, in general,  
$\tilde{s}$ does not satisfy conservation and $\tilde{e}$ does not satisfy compatibility.

We want to determine the best approximations to $\tilde{s}$ and $\tilde{e}$ among those fields that are feasible (i.e., conservative and compatible). The scalar field $u$ is obtained as a by-product.

\section{Well-posedness of the exact problem}

Let us define a few spaces:
\begin{eqnarray} 
U&=&\{\bar{u}\in H^1(\Omega), \bar{u}=0~\mbox{on}~\partial \Omega_D \}\\
E&=&L^2(\Omega)^d \\
S&=&L^2(\Omega)^d \\
X&=&U\times E \times S 
\end{eqnarray}

\bigskip

\noindent{\bf Problem P1:} Find $x=\left ( u,e,s\right )\in Z$ (the feasible space) that minimizes the distance to the data in the feasible space, i.e.,
\begin{equation}
J(x) \leq J(x+\delta x)~,\qquad \forall \,\delta x\,\in\,\delta Z~
\end{equation}
where 
\begin{equation}
J(\bar{u},\bar{e},\bar{s})=\frac12 \|\bar{e}-\tilde{e}\|^2 + \frac{\kappa}{2} \|\bar{s}-\tilde{s}\|^2~,
\end{equation}
the feasible space is
\begin{equation}
Z=\left \{ (\bar{u},\bar{e},\bar{s})\in X~|~\nabla \cdot \bar{s} + \zeta \,\bar{u} = q~, \bar{e}=\nabla \bar{u},\,\left. \bar{s}\cdot\mathbf{\check{n}}\right |_{\partial \Omega_N}=0 \right \}
\end{equation}
and its variation is
\begin{equation}
\delta Z=\left \{ (\delta{u},\delta{e},\delta{s})\in X~|~\nabla \cdot \delta{s} + \zeta \delta{u} = 0~, \delta{e}=\nabla \delta{u} ,\,\left. \delta{s}\cdot\mathbf{\check{n}}\right |_{\partial \Omega_N}=0\right \}~.
\end{equation}

\bigskip

Above $\|\cdot\|$ is the norm of $L^2(\Omega)^d$. The corresponding scalar product is denoted by $\left ( \cdot , \cdot \right )$. The same notation will be used for $L^2(\Omega)$.

Let us now relax the constraints through Lagrange multipliers $\lambda$ and $\mu$, belonging to 
$$
L=\{g\in H^1(\Omega)~|~g=0 \mbox{ on } \partial \Omega_D\},\qquad
\mbox{ and }\qquad M=L^2(\Omega)^d,  
$$
respectively. We will also use the notation
$$
y=\left (\lambda,\mu\right )\,\in \,Y\,=\,L\times M~.
$$
The resulting equivalent problem is

\bigskip

\noindent{\bf Problem P2:} Find $(x,y)=\left ( u, e, s, \lambda, \mu\right ) \subset X\times Y$ such that 
\begin{equation}
(x,y)= \mbox{arg} \left \{ \inf_{\bar{x}\in X} \sup_{\bar{y}\in Y} 
\mathcal{L}(\bar{x},\bar{y})
\right \}
\end{equation}
where
\begin{eqnarray}
\mathcal{L}(\bar{x},\bar{y})&=&J(\bar{u},\bar{e},\bar{s})+ \nonumber \\
& &-\left ( \bar{s},\nabla \bar{\lambda} \right ) + \left (\zeta \bar{u} - q , \bar{\lambda} \right )  
+ \left ( \bar{e}-\nabla \bar{u}, \bar{\mu} \right ) ~.
\end{eqnarray}

The weak form of the optimality conditions is

\begin{eqnarray}
 \left ( \lambda, \zeta\delta u \right ) - \left ( \mu, \nabla \delta u \right ) & = & 0 \qquad\qquad\qquad\qquad \forall \,\delta u\in U  \label{eq:var1}\\
\left ( e,\delta e \right )+ \left ( \mu, \delta e \right ) &=& \left ( \tilde{e},\delta e \right )  \qquad\qquad\qquad \forall \,\delta e \in E\label{eq:var2}\\
\kappa \left ( s, \delta s \right ) - \left ( \nabla \lambda, \delta s \right )  & = & \kappa \left ( \tilde{s},\delta s \right ) \qquad\qquad\qquad \forall\, \delta s \in S \label{eq:var3}\\
\left ( \zeta\,u,\delta \lambda \right )-\left ( s,\nabla \delta \lambda\right )  &=&\left ( q,\delta \lambda \right ) \qquad\qquad\qquad \forall\,\delta \lambda \in L\label{eq:var4}\\
-\left (\nabla u,\delta \mu \right )+\left ( e,\delta \mu \right ) & = & 0 \qquad\qquad\qquad\qquad \forall\,\delta \mu \in M \label{eq:var5}
\end{eqnarray}

\bigskip

The sum of the left-hand sides (multiplying the last two by $-1$) is the bilinear form of the problem in $X\times Y$,
\begin{equation}
B((x,y),(\delta x,\delta y))=a(x,\delta x)+b(\delta x,y)-b(x,\delta y)
\end{equation}
where
\begin{eqnarray}
a(x,\delta x)&=& \left (e,\delta e\right )+\kappa \left ( s,\delta s \right ) 
\end{eqnarray}
and
\begin{equation}
b(x,y)=\left ( \mu,e-\nabla u \right )-\left (\nabla \lambda,s\right ) + \left ( \lambda, \zeta u \right )
\label{eq:bform}
\end{equation}

\bigskip

The optimality conditions in strong form are
\begin{eqnarray}
\zeta \lambda  + \nabla \cdot \mu & = & 0 \label{eq:dif1}\\
 e + \mu & = & \tilde{e} \label{eq:dif2}\\
\kappa\, s - \nabla \lambda & = & \kappa\,\tilde{s} \label{eq:dif3}\\
 \zeta\,u+\nabla \cdot s  & = & q \label{eq:dif4}\\
-\nabla u + e &=&0 \label{eq:dif5}
\end{eqnarray}
together with
\begin{eqnarray}
 s \cdot \mathbf{\check{n}} & = & 0~,\qquad \mbox{ on } \partial \Omega_N \\
\mu\cdot\mathbf{\check{n}} & = & 0~,\qquad \mbox{ on } \partial \Omega_N ~.
\end{eqnarray}

\bigskip

\noindent{\bf Remark:} Notice that when the reaction coefficient $\zeta$ is identically zero the five variational equations (\ref{eq:var1})-(\ref{eq:var5})  (or, equivalently, their differential versions (\ref{eq:dif1})-(\ref{eq:dif5})) are in fact two independent sets of equations: (a) The first, second and fifth equations, with unknowns $u$, $e$ and $\mu$, and (b) the third and fourth equations, with unknowns $s$ and $\lambda$. In such a case, $s$ is the best approximation to $\tilde{s}$ among conservative flux fields, and $e=\nabla u$ is the best approximation to $\tilde{e}$ among compatible (gradient) fields, and both problems can be analyzed separately. The reaction term adds an interesting twist to the problem,  coupling all five variational equations.

\bigskip

%\noindent{\bf Brezzi's proof of convergence:}

\noindent{\bf Theorem 1:} Problem P2 has a unique solution.

\bigskip

\noindent{\em Proof:} Following the theory of mixed problems \cite{brezzifortin,ernguermond} one needs to show that $a(\cdot,\cdot)$ and $b(\cdot,\cdot)$ are continuous, that $a(\cdot,\cdot)$ is coercive in $\delta Z$, with
\begin{eqnarray}
\delta Z&=&\{x\in X~|~b(x,y)=0\,\forall y\in Y\}\nonumber \\
&=&\{(u,e,s)\in X|e-\nabla u=0, \nabla \cdot s + \zeta u =0,~\left. s\cdot \mathbf{\check{n}}\right |_{\partial\Omega_N} = 0  \}~,
\end{eqnarray}
and that the inf-sup condition on $b(\cdot,\cdot)$,
\begin{equation}
\inf_{y\in Y}\sup_{x\in X} \frac{b(x,y)}{\|x\|_X\|y\|_Y} \geq K_b > 0~.
\end{equation}
holds true.

The norms adopted in $X$ and $Y$ are 
\begin{eqnarray}
\|x\|_X &=& \|(u,e,s)\|_X = \sqrt{\|\nabla u\|^2+\|e\|^2+\kappa \|s\|^2}~,\\
\|y\|_Y &=& \|(\lambda,\mu)\|_Y=\sqrt{\frac{1}{\kappa}\left \| \nabla \lambda\right \|^2+\left\|\mu\right \|^2}~.
\end{eqnarray}
They were chosen so that they have the same units. It is easy to see that both $X$ and $Y$ are Hilbert spaces, and that $a(\cdot,\cdot)$ and $b(\cdot,\cdot)$ are continuous. Notice that
$$
a(x,x)=\|e\|^2 + \kappa \|s\|^2 ~.
$$
So, for $x\in \delta Z$ it holds that
$$
a(x,x)= \frac12 \|e\|^2 +\frac12 \|\nabla u\|^2+ \kappa \|s\|^2 \geq \frac12 \|x\|_X^2~.
$$
Coercivity thus holds with constant equal to $\frac12$.


It remains to show that there exists $K_b>0$ such that
\begin{equation}
\sup_{(u,e,s)\in X} \frac{\left (\mu,e-\nabla u \right ) - \left (\nabla \lambda, s \right ) + \left ( \lambda, \zeta u \right )}{\|(u,e,s)\|_X} \geq K_b \, \|(\lambda,\mu)\|_{Y}~,\qquad \forall (\lambda,\mu)\in Y~.
\end{equation}
We can take $u=0$, $e=\mu$
and $s=-\frac{1}{\kappa}\nabla \lambda$ to get
$$
\sup_{(u,e,s)\in X} \frac{\left (\mu,e-\nabla u \right ) - \left (\nabla \lambda, s \right ) + \left ( \lambda, \zeta\,u \right )}{\|(u,e,s)\|_X} 
\geq \frac{\|\mu\|^2 + \frac{1}{\kappa}\|\nabla \lambda\|^2}{\sqrt{\|\mu\|^2 + \frac{1}{\kappa}\|\nabla \lambda \|^2}} = \|(\lambda,\mu)\|_Y~,
$$
giving an inf-sup constant $K_b$ equal to one.

The coercivity of $a(\cdot,\cdot)$ on $\delta Z$ and the inf-sup stability of $b(\cdot,\cdot)$ have been proven. Then we have the well-posedness of the exact problem, with a unique solution $(u,e,s)\in H^1(\Omega)\times L^2(\Omega)^d \times L^2(\Omega)$ that depends continuously on $q\in L^2(\Omega)$, $\tilde{e}\in L^2(\Omega)^d$ and $\tilde{s}\in L^2(\Omega)^d$. $\Box$

\section{Approximation}

Let $U_h$, $E_h$, $S_h$, $L_h$ and $M_h$ be the finite element spaces for the five unknowns. We assume that $U_h\subset U$ and $L_h\subset L$ (in particular, the Dirichlet boundary conditions are strongly enforced). 
Also, we assume that $E_h$, $S_h$ and $M_h$ are subspaces of $L^2(\Omega)^d$ and define $X_h=U_h\times E_h \times S_h$ and $Y_h=L_h\times M_h$. The numerical approximation we propose corresponds to solving the following discrete problem:

\bigskip

\noindent{\bf Discrete Problem P2$_h$:} Find $(x_h,y_h)=\left ( u_h, e_h, s_h, \lambda_h, \mu_h\right ) \subset X_h\times Y_h$ such that 
\begin{equation}
(x_h,y_h)= \mbox{arg} \left \{ \inf_{\bar{x}\in X_h} \sup_{\bar{y}\in Y_h} 
\tilde{\mathcal{L}}(\bar{x},\bar{y})
\right \}
\end{equation}
where
\begin{eqnarray}
\tilde{\mathcal{L}}(\bar{x},\bar{y})&=&J(\bar{u},\bar{e},\bar{s})+ \nonumber \\
& &-\left ( \bar{s},\nabla \bar{\lambda} \right ) + \left (\zeta \bar{u} - q , \bar{\lambda} \right )  
+ \left ( \bar{e}-\nabla \bar{u}, \bar{\mu} \right ) +\nonumber \\
& &+\frac{\alpha}{2} \left \| \bar{e} - \nabla \bar{u} \right \|^2 + \\
& &-\frac{\gamma}{2} \|\bar{e}+\bar{\mu}-\tilde{e}\|^2 + \nonumber \\
& &-\frac{\eta}{2\kappa} \|\kappa \bar{s} - \nabla \bar{\lambda} - \kappa \tilde{s}\|^2 \nonumber \\
& &+ \frac{\theta \kappa \ell_s^2 }{2}\, \|\nabla \cdot \bar{s} + \zeta\, \bar{u} - q \|_h^2 + \nonumber \\
& & - \frac{\beta \ell_{\mu}^2 }{2}\, \|\nabla \cdot \bar{\mu} + \zeta\, \bar{\lambda} \|_h^2~.
\end{eqnarray}

\bigskip

The mesh-dependent norm $\|\cdot\|_h$ is defined as satisfying
\begin{equation}
\|f\|_h^2 = \sum_K \|f\|^2_{L^2(K)} 
\end{equation}
with the sum performed over all elements $K$ of the mesh. The associated scalar product is denoted as $\left ( \cdot, \cdot \right )_h$.

Above, $\alpha$, $\theta$ and $\beta$ are positive non-dimensional coefficients, and $\ell_s$ and $\ell_{\mu}$ are length scales to be determined later on. The function $\tilde{\mathcal{L}}$ is an augmented Lagrangian, of which the saddle point is the solution of P2$_h$. %%The solution is independent of $\alpha$ and $\theta$. 

The optimality conditions of P2$_h$ are

\begin{eqnarray}
\alpha \left ( \nabla u_h, \nabla \delta u \right ) + \theta \kappa \ell_s^2  \left ( \zeta^2 u_h, \delta u \right ) - \alpha \left ( e_h,\nabla \delta u \right ) +& & \nonumber \\
 + \theta \kappa \ell_f^2  \left ( \nabla \cdot s_h, \zeta \delta u \right )_h + \left ( \lambda_h, \zeta\delta u \right ) - \left ( \mu_h, \nabla \delta u \right ) & = & \theta \kappa \ell_s^2  \left ( q, \zeta\delta u \right )\nonumber \\ & & \quad \forall \,\delta u\in U_h  \label{eq:var1h}\\
-\alpha \left ( \nabla u_h, \delta e \right ) +(1+\alpha-\gamma)\,\left ( e_h,\delta e \right )+ (1-\gamma)\, \left ( \mu_h, \delta e \right ) &=& \left (1-\gamma)\,( \tilde{e},\delta e \right )  \nonumber \\ & &\quad \forall \,\delta e \in E_h\label{eq:var2h}\\
(1-\eta)\kappa \left ( s_h, \delta s \right ) - (1-\eta)\left ( \nabla \lambda_h, \delta s \right ) + & & \nonumber \\
+\theta \kappa \ell_s^2  \left ( \nabla \cdot s_h, \nabla \cdot \delta s \right )_h+ \theta \kappa \ell_s^2 \,\left (\zeta\,u_h, \nabla \cdot \delta s \right )_h & = & (1-\eta)\kappa \left ( \tilde{s},\delta s \right ) + \nonumber \\ & &+\theta \kappa \ell_s^2  \left ( q , \nabla \cdot \delta s \right )_h \nonumber \\ & & \quad \forall\, \delta s \in S_h \label{eq:var3h}\\
-(1-\eta)\left ( s_h,\nabla \delta \lambda\right ) + \left ( \zeta\,u_h,\delta \lambda \right ) - \beta \ell_{\mu}^2 \left ( \zeta^2 \lambda_h, \delta \lambda\right )+ & & \nonumber \\
- \beta \ell_{\mu}^2 \left ( \nabla \cdot \mu_h, \zeta\,\delta \lambda\right )_h -\frac{\eta}{\kappa} \left ( \nabla \lambda_h , \nabla \delta \lambda \right )&=&\left ( q,\delta \lambda \right ) + \eta \left ( \tilde{s},\nabla \delta \lambda \right ) \nonumber \\ & & \quad \forall\,\delta \lambda \in L_h\label{eq:var4h}\\
-\left (\nabla u_h,\delta \mu \right )+(1-\gamma)\left ( e_h,\delta \mu \right )
- \beta \ell_{\mu}^2 \left ( \zeta\,\lambda_h, \nabla \cdot \delta \mu \right )_h + & & \nonumber \\- \beta \ell_{\mu}^2 \left ( \nabla \cdot \mu_h, \nabla \cdot \delta \mu \right )_h - \gamma \left ( \mu_h,\delta \mu \right ) & = & -\gamma \left (\tilde{e},\delta \mu \right ) \nonumber \\ & & \quad \forall\,\delta \mu \in M_h \label{eq:var5h}
\end{eqnarray}

\bigskip

The sum of the left-hand sides (the last two multiplied by $-1$) is nothing but the discrete bilinear form that corresponds to P2$_h$, i.e.,
\begin{equation}
B_h((x_h,y_h),(\delta x,\delta y))=a_h(x_h,\delta x)+b(\delta x,y_h)-b(x_h,\delta y)+c_h(y_h,\delta y)~,
\end{equation}
where
\begin{eqnarray}
a_h(x_h,\delta x)&=&(1-\gamma)\left (e_h,\delta e\right )+(1-\eta)\kappa \left ( s_h,\delta s \right ) + \nonumber \\
& & +
\alpha \left ( e_h-\nabla u_h, \delta e - \nabla \delta u \right )+ \nonumber \\
& & + \theta \kappa \ell_s^2  \left ( \nabla \cdot s_h + \zeta u_h , \nabla \cdot \delta s + \zeta\,\delta u \right )_h~,
\end{eqnarray}
\begin{eqnarray}
b_h(\delta x,y_h)=b(\delta x,y_h)-\gamma \left (\mu_h,\delta e\right )+\eta \left (\nabla \lambda_h,\delta s \right ) ~,
\end{eqnarray}
\begin{eqnarray}
 c_h(y_h,\delta y)&= &  \beta \ell_{\mu}^2 \left ( \nabla \cdot \mu_h + \zeta \lambda_h , \nabla \cdot \delta \mu + \zeta\,\delta \lambda \right )_h+ \nonumber \\
 & & + \frac{\eta}{\kappa}\,\left ( \nabla \lambda_h,\nabla \delta \lambda \right ) + \gamma \left ( \mu_h,\delta \mu \right )~,
\end{eqnarray}
and $b(\cdot,\cdot)$ as defined in (\ref{eq:bform}).

Problem P2$_h$ can thus be written as "Find $(x_h,y_h)\in X_h \times Y_h$ such that
\begin{equation}
B_h\left ( \left ( x_h,y_h \right ),\left ( \delta x,\delta y \right ) \right )= F_h(\delta x,\delta y)
\label{eq:BzF}
\end{equation}
for all $(\delta x, \delta y)\in X_h\times Y_h$". The linear form $F_h:X_h\times Y_h \to \mathbb{R}$ stands for the sum of the right-hand sides of (\ref{eq:var1h})-(\ref{eq:var5h}), the last two multiplied by $-1$.

\bigskip

\noindent{\bf Theorem 2:} Problem P2$_h$ is well-posed, and $(x_h,y_h)$ satisfies that $\exists c>0$ independent of $h$ such that
\begin{equation}
\||(x,y)-(x_h,y_h)\||_h \leq c\,\min_{(\bar{x},\bar{y})\in X_h\times Y_h}
\||(x,y)-(\bar{x},\bar{y})\||_h
\label{eq:bestapprox}
\end{equation}
where
\begin{eqnarray}
\||(x_h,y_h)\||_h^2 &=& \|\nabla u_h \|^2 +\|e_h\|^2 + \kappa \|s_h\|^2 + \kappa h^2 \|\nabla \cdot s_h\|_h^2 + \nonumber \\
& &+ \frac{1}{\kappa} \|\nabla \lambda_h\|^2 + \|\mu_h\|^2 + h^2\,\|\nabla \cdot \mu_h\|^2_h~.
\end{eqnarray}

\bigskip

The terms added to the Lagrangian allow us to work in a norm with improved control over $\nabla u_h$, $\nabla \cdot s_h$ and $\nabla \cdot \mu_h$.

\bigskip

\noindent{\em Proof:} It is easy to verify that the exact solution $(x,y)$ satisfies the discrete problem (\ref{eq:BzF}) when substituted in the place of $(x_h,y_h)$. The formulation is thus consistent.

 Let us establish the uniform strong coercivity of $B_h$, i.e., the existence of $k_{B}>0$ (independent of $h$) such that
\begin{equation}
B_h((x_h,y_h),(x_h,y_h))\geq k_B \left ( \|x_h\|_X^2 + \|y_h\|_Y^2 \right ),\qquad \forall (x_h,y_h)\in X_h\times Y_h~.
\end{equation}
By direct computation, 
\begin{eqnarray}
B_h((x_h,y_h),(x_h,y_h))&=&(1-\gamma) \|e_h\|^2 + (1-\eta)\kappa \|s_h\|^2 + \nonumber \\
& & + \alpha \|e_h-\nabla u_h\|^2 + \theta \kappa \ell_s^2 \|\nabla \cdot s_h + \zeta u_h\|^2 + \nonumber \\
& & + \beta \ell_{\mu}^2 \| \nabla \cdot \mu_h + \zeta \lambda_h \|^2
+ \frac{\eta}{\kappa} \|\nabla \lambda_h\|^2 + \gamma \|\mu_h\|^2~.
\end{eqnarray}


Now we use Young's inequality ($2ab\geq -\epsilon^2 a^2-(1/\epsilon^2)b^2$, $\forall a,b, \epsilon$) to each of the cross-product terms in the previous expression:

\begin{eqnarray}
-2\alpha \left (e_h,\nabla u_h \right ) & \geq & -\frac14 \|e_h\|^2 - 4\alpha^2 \|\nabla u_h \|^2 \\
%%%%%
2\theta\kappa \ell_s^2\left ( \nabla \cdot s_h, \zeta u_h \right )_h& \geq & - \frac12 \kappa \theta \ell_s^2 \|\nabla\cdot s_h\|^2 - {2\theta\kappa\ell_s^2}\|\zeta u_h\|^2\\
%%%%%%
2 \beta \ell_{\mu}^2\left ( \nabla \cdot \mu_h, \zeta \lambda_h \right )_h& \geq & - \frac12 {\beta} \ell_{\mu}^2 \|\nabla\cdot \mu_h\|^2 - {2\beta\ell_{\mu}^2}\|\zeta \lambda_h\|^2\\
%%%%%%%
\end{eqnarray}

Adding up all terms that contain $s_h$ we arrive at
\begin{equation}
(1-\eta) \kappa \|s_h\|^2 + \frac12 \kappa \ell_s^2 
 \theta  \|\nabla\cdot s_h\|_h^2~.
\label{eq:term1}
\end{equation}
Adding up all terms that contain $\zeta u_h$ gives
$$
-\theta \kappa \ell_s^2 \|\zeta u_h\|^2
$$
which using Poincar\'e-Friedrichs inequality
\begin{equation}
\|u_h\| \leq c_{\mbox{\scriptsize{PF}}}\,\ell_{\Omega}\,\|\nabla u_h\|, \qquad \forall u_h\in U_h
\label{eq:poincare}
\end{equation}
with $\ell_{\Omega}$ the diameter of $\Omega$ is shown to be greater or equal than 
\begin{equation}
- {\theta \kappa \ell_s^2 \|\zeta\|_{\infty}^2 c_{\mbox{\scriptsize{PF}}}^2 \ell_{\Omega}^2}\,\|\nabla u_h \|^2~.
\label{eq:term2}
\end{equation}
By $\|\zeta\|_{\infty}$ we denote the $L^{\infty}(\Omega)$-norm of $\zeta$.

Adding up all terms that contain $\mu_h$ we arrive at
\begin{equation}
\gamma \|\mu_h\|^2 +\frac12 \beta\ell_{\mu}^2  \|\nabla \cdot \mu_h \|_h^2~.
\label{eq:term4}
\end{equation}
The sum of the terms that contain $\zeta \lambda_h$ is greater or equal, by Poincaré-Friedrichs inequality, than
\begin{equation}
- \beta \ell_{\mu}^2 \|\zeta\|_{\infty}^2 c_{\mbox{\scriptsize{PF}}}^2 \ell_{\Omega}^2\,\|\nabla \lambda_h \|^2~.
\label{eq:term5}
\end{equation}
Adding up all terms that contain $\nabla u_h$ gives
\begin{equation}
\left ( \alpha - 4 \alpha^2 \right )
\|\nabla u_h\|^2~.
\label{eq:term6}
\end{equation}
Adding up all terms that contain $e_h$ gives
\begin{equation}
\left ( \frac34-\gamma + \alpha \right ) \|e_h\|^2~.
\label{eq:term7}
\end{equation}

Adding everything up we get
\begin{eqnarray}
B_h((x_h,y_h),(x_h,y_h))&\geq& \left (\frac34 -\gamma+\alpha \right ) \|e_h\|^2 + (1-\eta)\kappa \|s_h\|^2 + \nonumber \\
& & + (\alpha - 4 \alpha^2 - \theta \kappa \ell_s^2 \|\zeta\|_{\infty}^2 c_{\mbox{\scriptsize{PF}}}^2 \ell_{\Omega}^2)\, \|\nabla u_h\|^2 + \nonumber \\
& & + \frac12 \theta \kappa \ell_s^2 \|\nabla \cdot s_h\|^2 + \nonumber \\
& & + \frac12 \beta \ell_{\mu}^2 \| \nabla \cdot \mu_h\|^2 + \nonumber \\
& & 
+ \left ( \frac{\eta}{\kappa}-\beta \ell_{\mu}^2 \|\zeta\|_{\infty}^2 c_{\mbox{\scriptsize{PF}}}^2 \ell_{\Omega}^2 \right ) \|\nabla \lambda_h\|^2 + \gamma \|\mu_h\|^2~.
\end{eqnarray}

Taking
$$
\alpha = \frac{1}{8}, \quad \theta=\frac12,\quad \beta=\frac12,
\quad \gamma=\alpha, \quad \eta = \frac12,
\quad \ell_s = \ell_\mu = h, 
$$
we get
\begin{eqnarray*}
B_h\left ( (x_h,y_h),(x_h,y_h) \right ) &\geq  &
\left ( \frac{1}{16} - O(h^2) \right )\, \|\nabla u_h\|^2+\\
& & +\frac{3}{4}\,\|e_h\|^2+\frac12 \kappa \|s_h\|^2 + \frac14 \kappa h^2 \|\nabla \cdot s_h\|_h^2 + \\
& &
+\left (\frac{1}{2\kappa}  -O(h^2) \right ) \|\nabla \lambda_h\|^2 + \\
& & + \frac18 \|\mu_h\|^2  +\frac14 h^2 \|\nabla \cdot \mu_h\|_h^2 
\end{eqnarray*}
and finally
\begin{equation}
B_h\left ( (x_h,y_h),({x}_h,{y}_h) \right ) \geq 
k_B \, \| |(x_h,y_h)  \| |_h^2 
\end{equation}
where
\begin{equation}
k_B \geq \frac{1}{16}~.
\end{equation}
We have assumed that $h$ is sufficiently small so that the terms $O(h^2)$ can be neglected and thus $k_B$ becomes mesh independent. 

It is straightforward to see that $B_h(\cdot,\cdot)$ is continuous, i.e., $\exists N_B\geq 0$ such that
\begin{equation}
B_h\left ( (x_h,y_h),(\bar{x},\bar{y}) \right ) \leq N_B \||(x_h,y_h)\||_h\,\||(\bar{x},\bar{y})\||_h~,
\end{equation}
$ \forall (x_h,y_h) \in X_h\times Y_h$, $\forall (\bar{x},\bar{y})\in X_h\times Y_h$, with $N_B$ independent of $h$. Having proved that $B_h$ is consistent, continuous and coercive we obtain the best approximation property (\ref{eq:bestapprox}). $\Box$


\end{document}
